% REVISION: all tables below retain the supplied old measurements. Replace from corrected raw logs, not ratios alone.
\section{Experimental Results}
\label{sec:experimental-results}
This section evaluates the performance of our AF-CPSI and the resulting
PJC protocol that combines it with homomorphic-encryption-based
computation. We first describe the implementation environment and
cryptographic parameters, then measure the running time and communication
of our protocol across a range of input sizes and network settings and
compare it against existing PJC protocols. Finally, we break down the
protocol's per-stage cost to identify the main contributors to overall
performance.



\subsection{Implementation and Experimental Setup}
\label{sec:experimental-setup}

\paragraph{Experiment environments.}
All experiments were run on Ubuntu 24.04.4 LTS with an AMD Ryzen
Threadripper 9970X 32-Core CPU and 64\,GB of RAM. Each protocol was
executed on a single thread. We evaluated two network settings: LAN and
emulated WAN. The LAN setting used localhost with no bandwidth
restriction, and the WAN setting used Linux \textsf{tc} to cap bandwidth
at 100\,Mbps and 10\,Mbps, respectively. We added no additional network
latency, and confirmed that no packet drops occurred in any of the WAN
measurements.

\medskip
\paragraph{Implementation details.}
We implemented our protocol in C++. AF-CPSI builds on \textsf{VOLE-PSI}~\cite{volepsi},
and we use \textsf{libOTe}~\cite{libOTe} for OT and B2A-OT, and we implement the homomorphic-encryption-based PJC
computation using the BGV~\cite{BGV14} scheme provided by Microsoft SEAL~\cite{SEAL}.

\medskip
\paragraph{Parameters.}
We set all cryptographic parameters to satisfy 128-bit computational
security and 40-bit statistical security. For BGV, we use polynomial
modulus degree \(N=4096\), plaintext modulus \(p\approx 2^{32}\), and
ciphertext modulus \(q\approx 2^{84}\); this \(q\) suffices since
homomorphic evaluation consists of a single ciphertext--plaintext
multiplication followed by modulus switching. Batching packs up to
\(N=4096\) values into a single ciphertext. In the integer inner product
setting (\Cref{sec:integer-ip}) we use two 42-bit primes as the plaintext
moduli (\(p_1,p_2\approx 2^{42}\)) and a ciphertext modulus
\(q\approx 2^{105}\). Table~\ref{tab:experimental-parameters} summarizes
the parameters used in our experiments.


\begin{table}[!htbp]
\centering
\caption{Cryptographic parameters used in our experiments.}
\label{tab:experimental-parameters}
\begin{tabular}{ll}
\toprule
\textbf{Parameter} & \textbf{Setting} \\
\midrule
HE scheme & BGV \\
Polynomial modulus degree $N$ & \(4096\) \\
Plaintext modulus $p$ & \(\approx 2^{32}\) \\
Ciphertext modulus $q$ & \(\approx 2^{84}\) \\
\midrule
\multicolumn{2}{l}{\textit{Integer inner product setting (\Cref{sec:integer-ip})}} \\
Plaintext moduli (RNS) & \(p_1,p_2\approx 2^{42}\) \\
Ciphertext modulus $q$ & \(\approx 2^{105}\) \\
\bottomrule
\end{tabular}
\end{table}


\subsection{Comparison: Modular Inner Product}

\label{sec:performance-comparison}
This section compares the running time and communication of PJC
protocols across different input sizes and network settings. We compare
\cite{CHI+23}, \cite{KLS26}, and our protocol. \cite{KLS26} and our
protocol are evaluated in the same experimental environment under
$\ell=L=32$, i.e., 32-bit payloads and a 32-bit arithmetic space.
\cite{CHI+23}, in contrast, has no public implementation, so we report
the numbers as given in the original paper. Since \cite{CHI+23} uses
$\ell=23$, the communication comparison is favorable to
\cite{CHI+23}.



\begin{table}[!htbp]
\centering
\caption{Performance comparison of PJC protocols under different network settings (32-bit setting).}
\label{tab:pjc-comparison}

\setlength{\tabcolsep}{3.5pt}
\renewcommand{\arraystretch}{1.05}
\small

\begin{tabular}{c l c c c c}
\toprule

\multirow{2}{*}{$n$}
&
\multirow{2}{*}{\textbf{Protocol}}
&
\multirow{2}{*}{\textbf{Comm. (MB)}}
&
\multicolumn{3}{c}{\textbf{Runtime (s)}}
\\

\cmidrule(lr){4-6}

&
&
&
\textbf{LAN}
&
\textbf{WAN}
&
\textbf{WAN}
\\[-2pt]

&
&
&
&
\textbf{100 Mbps}
&
\textbf{10 Mbps}
\\

\midrule

\multirow{3}{*}{$2^{16}$}
& \cite{CHI+23}$^\dagger$& 17.2 & 25.92 & -- & -- \\
& \cite{KLS26}&20.90&1.74&3.28&18.15\\
\cmidrule(lr){2-6}%
& \textbf{Ours}&14.15&0.74&1.56&11.38\\

\midrule

\multirow{3}{*}{$2^{20}$}
& \cite{CHI+23}$^\dagger$& 275.2 & 424.4 & -- & -- \\
& \cite{KLS26}&322.62&28.11&53.53&287.51\\
\cmidrule(lr){2-6}
& \textbf{Ours}&218.01&11.39&26.84&182.45\\

\bottomrule
\end{tabular}

\vspace{2pt}
\begin{minipage}{0.95\linewidth}
\footnotesize
$^\dagger$ \cite{CHI+23} has no public implementation, so we report the LAN numbers from the original paper.
\end{minipage}

\end{table}

As shown in Table~\ref{tab:pjc-comparison}, our protocol achieves lower
running time and communication than \cite{KLS26} across all input sizes
and network settings. At set size $2^{16}$, our protocol's communication
cost is 14.15\,MB, about $1.48\times$ lower than \cite{KLS26}'s
20.90\,MB, and its LAN running time is likewise about $2.35\times$
faster, at 0.74\,s versus 1.74\,s.
 
This trend holds at set size $2^{20}$ as well. Our protocol's
communication cost of 218.01\,MB is about $1.48\times$ lower than
\cite{KLS26}'s 322.62\,MB, and its LAN running time of 11.39\,s is about
$2.47\times$ faster than 28.11\,s. Under 100\,Mbps and 10\,Mbps WAN
settings, our protocol runs in 26.84\,s and 182.45\,s, respectively, and
continues to outperform \cite{KLS26}. This shows that our protocol's
lower communication cost translates into a running-time advantage even
under constrained bandwidth.
 
Since \cite{CHI+23} was evaluated under a 23-bit setting, its
communication cost in Table~\ref{tab:pjc-comparison} is not measured
under the same conditions as the 32-bit setting used for \cite{KLS26}
and our protocol; we include it only as a reference point for the
communication cost of prior PJC protocols. We report a measured
comparison for the integer-inner-product setting over 32-bit payloads in
\Cref{sec:integer-ip}.


\paragraph{Breakdown.}

We measured the per-stage running time and communication
over LAN for set size $2^{20}$, with three stages
-- AF-CPSI, B2A-OT, and homomorphic-encryption-based PJC computation --
and report the results in Table~\ref{tab:pjc-breakdown}.

\begin{table}[!htbp]
\centering
\caption{Breakdown of our PJC protocol over LAN for $n=2^{20}$.}
\label{tab:pjc-breakdown}
\begin{tabular}{lcc}
\toprule
\textbf{Stage}
&
\textbf{Runtime (s)}
&
\textbf{Comm. (MB)}
\\
\midrule
AF-CPSI&10.74&172.96\\
B2A-OT&0.36&22.77\\
HE Computation&0.26&22.29\\
\midrule
\textbf{Total}&11.39&218.01\\
\bottomrule
\end{tabular}
\end{table}

As Table~\ref{tab:pjc-breakdown} shows, the bulk of the total cost is
incurred during AF-CPSI, which dominates both running time and
communication and stands out as the main performance bottleneck of our
protocol. 

\subsection{Comparison: Integer Inner Product}
\label{sec:integer-ip-exp}
We measured the setting of \Cref{sec:integer-ip}, i.e., the exact integer
inner product over 32-bit payloads, under LAN and two WAN settings.
 
The implementation of \cite{KLS26} is fixed to $\Z_{2^{32}}$ and cannot
be measured under a different ring without modification, so we measured
it using a cost harness with all OKVS value, OT message, and MultShare
widths matched to 84 bits. This harness preserves every width that
determines cost in the original protocol, but replaces the underlying
operations with XOR and thus does not produce an inner-product value.

\begin{table}[!htbp]
\centering
\caption{Integer inner product over 32-bit payloads ($L=84$) under different network settings.}
\label{tab:pjc-long}
\setlength{\tabcolsep}{3.4pt}
\renewcommand{\arraystretch}{1.05}
\small
\begin{tabular}{c l c c c c}
\toprule
\multirow{2}{*}{$n$} & \multirow{2}{*}{\textbf{Protocol}} & \multirow{2}{*}{\textbf{Comm. (MB)}} & \multicolumn{3}{c}{\textbf{Runtime (s)}} \\
\cmidrule(lr){4-6}
& & & \textbf{LAN} & \textbf{100 Mbps} & \textbf{10 Mbps} \\
\midrule
\multirow{2}{*}{$2^{16}$}
& \cite{KLS26} & 93.53 & 2.68 & 9.93 & 77.34 \\
& \textbf{Ours} & 21.83 & 0.82 & 2.24 & 17.64 \\
\midrule
\multirow{2}{*}{$2^{20}$}
& \cite{KLS26} & 1499.29 & 43.05 & 162.30 & 1255.18 \\
& \textbf{Ours} & 340.63 & 13.39 & 38.32 & 282.59 \\
\bottomrule
\end{tabular}
\vspace{2pt}
\begin{minipage}{0.95\linewidth}
\footnotesize
\cite{KLS26} is measured with a width-matched cost harness (see text).
\cite{CHI+23} has no public implementation and reports only $\ell=23$.
\end{minipage}
\end{table}

The results in Table~\ref{tab:pjc-long} relate to the 32-bit setting of
Table~\ref{tab:pjc-breakdown} as follows. With two residues, the OKVS
data component, B2A-OT, and HE stages each double, while the identifier
path is unchanged. At $n=2^{20}$, our protocol's communication grows
from 218.01\,MB to 340.63\,MB, about $1.56\times$, whereas
\cite{KLS26}'s grows from 322.62\,MB to 1499.29\,MB, about $4.65\times$,
owing to the quadratic growth.
 
Head to head in this setting, our protocol uses $4.40\times$ less
communication at $n=2^{20}$ and $4.28\times$ less at $n=2^{16}$. The LAN
running times are shorter by $3.22\times$ and $3.27\times$; the LAN link is
fast enough that computation, rather than transmission, sets the pace, so
the running-time gap there falls short of the communication gap. The two
converge as bandwidth drops. At $n=2^{20}$ our protocol is $4.24\times$
faster over 100\,Mbps and $4.44\times$ faster over 10\,Mbps, at which point
the running-time gap has caught up with the communication gap.
